\documentclass[11pt]{article}
\usepackage{amsmath,amssymb,amsthm,mathtools}
\usepackage[margin=1in]{geometry}
\usepackage{multirow}
\usepackage{hyperref}
\newtheorem{theorem}{Theorem}
\newcommand{\Q}{\mathbb{Q}}\newcommand{\Z}{\mathbb{Z}}
\newcommand{\zetaK}{\zeta_{K}}\newcommand{\F}{\mathbb{F}}

\title{The Positivity Wall: Geometric Routes to the\\ Riemann Hypothesis, and Where
They Stop}
\author{(VFD programme --- expository working draft)}
\date{2026}

\begin{document}
\maketitle

\begin{abstract}
This is an \emph{expository} synthesis. We recall how the Riemann Hypothesis (RH)
reduces to a single positivity, survey the many equivalent forms this positivity
takes (spectral, algebraic, dynamical, geometric), and explain why every geometric
``route'' we examined reaches the same inequality without crossing it. We then state,
as precisely as the literature allows, the single object whose construction would
close it---a canonical Frobenius over $\Z$ whose degree-form is the Weil positivity,
i.e.\ the arithmetic surface of the function-field analogue. \textbf{No new result on
RH is claimed.} Every mathematical statement below is standard and credited; the
contribution is organizational: a clean map of the boundary, with explicit stamps
separating what is known, what is verified here, and what is interpretation.
\end{abstract}

\section{Scope}
We claim no theorem about RH. The value of this note is a single, honestly stamped
picture: RH is one inequality; it appears in many costumes; geometry stages it but
does not close it; and the missing step is a construction, not a search. Readers
seeking the verified arithmetic of the present programme are referred to the
companion papers; here we only organize known mathematics.

\section{RH as one positivity}
By the Weil explicit-formula criterion \cite{Weil}, RH is equivalent to the
nonnegativity of an explicit quadratic functional $W(f)\ge 0$ on a suitable space of
test functions. Every route below is a re-expression of this one inequality.

\section{One condition, many faces}
The following are standard equivalent criteria and standard programme templates for
RH. The first three are \emph{proven equivalent} to RH; the last two are programme
templates (one open over $\Z$, one a theorem in the function-field model). A single
off-line zero falsifies the proven-equivalent criteria simultaneously:
\begin{itemize}\itemsep1pt
  \item Weil positivity $W(f)\ge 0$ (spectral / Connes \cite{Connes}) --- equivalent
        to RH;
  \item the Li coefficients $\lambda_n\ge 0$ (algebraic \cite{Li}) --- equivalent to
        RH;
  \item the de Bruijn--Newman constant $\Lambda\le 0$ (dynamical; with
        $\Lambda\ge 0$ proven \cite{RodgersTao}, so RH $\iff \Lambda=0$);
  \item hyperbolicity of \emph{all} Jensen polynomials of $\Xi$ is equivalent to RH;
        GORZ \cite{GORZ} prove the fixed-degree/asymptotic part, not RH itself;
  \item the existence of a self-adjoint operator with the zeros as spectrum
        (Hilbert--P\'olya) --- a \emph{programme}, not a proven equivalence;
  \item positivity of a Frobenius degree-form on a surface --- a \emph{theorem} in
        the function-field template \cite{Weil2,Deligne}, with no built analogue over
        $\Z$.
\end{itemize}
Viewed as the proven-equivalent core, these are not many problems but one, seen from
many sides; the last two are the geometric/spectral templates one hopes to transport.

\section{The structure around the zeros (true, but not the wall)}
The Riemann zeros are modelled, with strong numerical and partial-theorem support, by
the statistics of a chaotic, time-reversal-broken (GUE, $\beta=2$) spectrum: Montgomery
\cite{Montgomery} proves the pair correlation in a restricted range (conditionally),
and Odlyzko \cite{Odlyzko} confirms it numerically to high precision; full GUE
statistics remain conjectural. Summed as oscillations the zeros reproduce the primes
with von-Mangoldt weights (the explicit formula, unconditional); and Rodgers--Tao give
$\Lambda\ge 0$, so RH would force the de Bruijn--Newman constant to $\Lambda=0$ (not
known unconditionally). The explicit formula holds whether or not RH is true; the GUE
statistics are a separate, largely empirical input. These constrain the shape of any
solution; they do not supply the positivity.

\section{The missing object}
The function-field analogue of RH is a theorem precisely because the positivity is
\emph{geometric}: for a curve $C/\F_q$, the bound on Frobenius eigenvalues follows
from the Hodge-index theorem on the surface $C\times C$. Over $\Z$ the analogous
object---a canonical Frobenius whose degree-form is $W(f)$, living on a putative
arithmetic surface $\operatorname{Spec}\Z\times_{\F_1}\operatorname{Spec}\Z$---is
\emph{unbuilt}. The leading programme towards it---Connes' trace-formula /
noncommutative-geometry approach and the related $\F_1$ / absolute-geometry
programmes (Connes--Consani \cite{ConnesConsani})---recovers the explicit formula as a
trace, but reduces RH to a \emph{positivity of a trace pairing} that remains unproven.
(Brave-new-ring heuristics model $\F_1$ by the sphere spectrum $\mathbb{S}$ and the
surface by $\Z\otimes_{\mathbb{S}}\Z$; this is a separate suggestion, not part of
Connes--Consani's site.) The positivity is the wall.

\section{A table of symbols, and which one is a placeholder}
The RH sentence can be written entirely in standard notation; the notation is
\emph{complete}. What is incomplete is the \emph{referent} of one symbol: it names a
hoped-for object that has no agreed construction. We make this explicit.

\begin{center}\small
\begin{tabular}{lll}
\hline
symbol & reading & status \\
\hline
$\operatorname{Re}(s)=\tfrac12$ & critical line & defined; computable \\
$|\alpha_i|=q^{1/2}$ & Frobenius eigenvalues on a circle & \emph{theorem} (function field) \\
$\xi(s)=\xi(1-s)$ & the mirror / functional equation & defined; proven \\
$W(f)\ge 0$ & Weil positivity & defined; \emph{unproven} ($=$ RH) \\
$\operatorname{Frob}_q$ & geometric Frobenius & a real morphism \\
$\operatorname{Spec}\Z\times_{\F_1}\operatorname{Spec}\Z$ & arithmetic surface & \textbf{no construction carrying} \\
 & & \textbf{the required positivity} \\
$\F_1$ & ``field with one element'' & \textbf{many frameworks; none} \\
 & & \textbf{supplies the RH object} \\
\hline
\end{tabular}
\end{center}

\noindent Several inequivalent $\F_1$-frameworks exist (\S\ref{sec:ledger}); what none
supplies is a canonical arithmetic surface carrying the required intersection theory
and positivity. The reciprocal mirror $x\mapsto 1/x$ on $\mathbb{P}^1$ (radius $0
\leftrightarrow \infty$, the unit circle $|x|=1$ fixed) is a schematic model of the
involution behind $\xi(s)=\xi(1-s)$ under $s=\tfrac12+it \leftrightarrow \log|x|$; that
involution is built. The unbuilt entry is the $\F_1$-surface. Writing the symbol is
legitimate; \emph{constructing the object it should point to is exactly the open
problem.}

\section{A ledger of \texorpdfstring{$\F_1$}{F1} candidates}\label{sec:ledger}
Because $\F_1$ is the placeholder, it is natural to ask whether one may simply
\emph{define} it---``$\F_1 :=$ unity,'' say---and proceed. Many definitions exist; the
question is never whether it can be \emph{named} but whether the named object
\emph{carries} the explicit-formula positivity \emph{without assuming it}. We classify
the principal published candidates by three documented criteria: \textsc{defined} (an
agreed object exists), \textsc{carries} (building the surface over it yields $W(f)$ as a
genuine degree/intersection form), and \textsc{honest} (it does so without putting the
conclusion in by hand). The verdicts below report the literature's status, not a
computation: whether a given $\F_1$ proves RH \emph{is} the open problem.

\begin{center}\small
\begin{tabular}{lll}
\hline
candidate & verdict & where it stops \\
\hline
$\F_1=\{*\}$, trivial monoid (``unity'') & \textsc{inert} & no nontrivial Frobenius \\
$\F_1$-vector space $=$ pointed set \cite{Cohn} & \textsc{inert} & $q\to1$ kills the arithmetic \\
Deitmar monoid schemes \cite{Deitmar} & \textsc{inert} & toric-like (finite-type) factors \\
Borger $\Lambda$-rings \cite{Borger} & \textsc{open} & supplies Frobenius lifts, not $W(f)$ \\
Connes--Consani arithmetic site \cite{ConnesConsani} & \textsc{open} & trace recovered, positivity open \\
Deninger dynamical cohomology \cite{Deninger} & \textsc{conjectural} & cohomology hypothetical \\
``$\F_1$ defined so $|\alpha|=q^{1/2}$'' (control) & \textsc{circular} & assumes the conclusion \\
\hline
\end{tabular}
\end{center}

\noindent \emph{None of the listed candidates is known to supply the required
positivity without an extra hypothesis.} Every honestly-defined object is either
\textsc{inert} (real, but carries no zeros---the fate of ``unity''), \textsc{open}
(Connes--Consani, Borger: reaches the wall, positivity not established), or
\textsc{conjectural} (Deninger: the needed cohomology is hypothetical). The only
candidate that \emph{carries} the positivity outright is the \textsc{circular} control,
which assumes $|\alpha|=q^{1/2}$ and is rejected on the same provenance grounds that
reject a least-squares fitted map. The honest summary: one can name a real object, or
one that would cross; no listed framework is known to be both.

\section{The joint void: a single door seen from two rooms (a conjectural picture)}
\emph{This section is an interpretive picture, not a theorem.} RH is approached from two
sides that the Hilbert--P\'olya philosophy and Connes' programme \emph{hope} are one
object seen twice: an \emph{arithmetic} side (a Frobenius on the $\F_1$-surface whose
degree-form would be $W(f)\ge 0$) and a \emph{spectral} side (a self-adjoint $H=H^*$
whose spectrum would be the zeros---Hilbert--P\'olya, Berry--Keating \cite{BerryKeating}).
Tabulating what each side reaches today (``reach'' = a partial construction exists;
``part.'' = modelled but not fully built; ``dark'' = no construction):

\begin{center}\small
\begin{tabular}{llll}
\hline
requirement & arithmetic & spectral & \\
\hline
explicit/trace formula & reach & reach & \multirow{4}{*}{\small shared floor (no RH)}\\
global density (Weyl law) & reach & reach & \\
GUE local statistics (empirical) & part. & part. & \\
mirror $\xi(s)=\xi(1-s)$ & reach & reach & \\
\hline
Frobenius / arithmetic surface & part. & \emph{dark} & geometry's half (unbuilt) \\
self-adjoint $H$ on a positive space & \emph{dark} & part. & physics' half (no final $H$) \\
\hline
$W(f)\ge0$ \ /\ ($\operatorname{spec} H=$ zeros, $H=H^*$) & \emph{dark} & \emph{dark} & \textbf{the door} \\
\hline
\end{tabular}
\end{center}

\noindent Two observations. First, the two sides \emph{already meet} on the structure
that does not require RH (the shared floor; GUE only empirically). Second, the
programme's \emph{hope} is that the arithmetic-dark and spectral-dark requirements are
the same object missing its two faces. We emphasise the status carefully: Connes'
analysis shows RH is equivalent to the positivity of a \emph{trace pairing}; it does
\emph{not} prove that ``$W(f)\ge 0$'' and ``$H=H^*$'' are literally the same statement,
and \textbf{no theorem identifies the arithmetic and spectral constructions.} What can
be said honestly is that \emph{both, under their own strong hypotheses, would yield RH},
and the programme conjectures a single underlying object. If that conjecture held, the
``both'' requirement would not widen the problem but pin the missing object toward one
shape---a heuristic motivation, not a lever. It does not lower the wall: supplying the
positivity by fiat on either face is circular, as above. The honest gain is only a
sharper \emph{picture}---``a single door, seen from two rooms''---with the floor on both
sides and each doorpost, at most, half-built.

\section{Closing}
Every accessible route---cavity, mirror, operator, loop, gas, flow, hyperbolicity---is
a re-expression of one positivity, and re-expression does not prove it. The single
missing step is the construction of an arithmetic surface and the positivity of its
intersection form. We make no claim of progress on that step; we have only drawn the
boundary cleanly, and marked, at each face, exactly where known mathematics ends.

\paragraph{A note on interpretation.} Several physically suggestive readings arise
naturally when one draws these pictures (curvature, mirrors, resonance). They are
\emph{interpretation only}; no statement in this note depends on them, and none is
claimed as mathematics.

\begin{thebibliography}{99}
\bibitem{Weil} A.~Weil, \emph{Sur les ``formules explicites'' de la th\'eorie des
nombres premiers}.
\bibitem{Connes} A.~Connes, \emph{Trace formula in noncommutative geometry and the
zeros of the Riemann zeta function}, Selecta Math.\ (1999).
\bibitem{Li} X.-J.~Li, \emph{The positivity of a sequence of numbers and the Riemann
Hypothesis}, J.\ Number Theory (1997).
\bibitem{GORZ} M.~Griffin, K.~Ono, L.~Rolen, D.~Zagier, \emph{Jensen polynomials for
the Riemann zeta function and other sequences}, PNAS (2019).
\bibitem{RodgersTao} B.~Rodgers, T.~Tao, \emph{The de Bruijn--Newman constant is
non-negative}, Forum Math.\ Pi (2020).
\bibitem{Weil2} A.~Weil, \emph{Sur les courbes alg\'ebriques et les vari\'et\'es qui
s'en d\'eduisent} (1948).
\bibitem{Deligne} P.~Deligne, \emph{La conjecture de Weil I}, IHES (1974).
\bibitem{Montgomery} H.~Montgomery, \emph{The pair correlation of zeros of the zeta
function} (1973).
\bibitem{Odlyzko} A.~Odlyzko, \emph{On the distribution of spacings between zeros of
the zeta function} (1987).
\bibitem{ConnesConsani} A.~Connes, C.~Consani, \emph{Schemes over $\F_1$ and zeta
functions}; and subsequent work on the arithmetic/scaling site.
\bibitem{Cohn} H.~Cohn, \emph{Projective geometry over $\F_1$ and the Gaussian
binomial coefficients}, Amer.\ Math.\ Monthly (2004).
\bibitem{Deitmar} A.~Deitmar, \emph{Schemes over $\F_1$}, in \emph{Number Fields and
Function Fields} (2005); see also B.~To\"en, M.~Vaqui\'e, \emph{Au-dessous de
$\operatorname{Spec}\Z$} (2009).
\bibitem{Borger} J.~Borger, \emph{$\Lambda$-rings and the field with one element}
(2009).
\bibitem{Deninger} C.~Deninger, \emph{Some analogies between number theory and
dynamical systems on foliated spaces}, ICM (1998).
\bibitem{BerryKeating} M.~V.~Berry, J.~P.~Keating, \emph{The Riemann zeros and
eigenvalue asymptotics}, SIAM Review (1999).
\end{thebibliography}

\end{document}
